% ============================================================================
%  c_tmd snippet 4/4: Sivers 函数的过程依赖（T11.14）与单胶子末态相互作用相位（E11.16）
%  符号：Trento 约定 f = f_1 - f_1T^perp (P^ x k_T).S / M；Dirac 表象，gamma_5 = i g0 g1 g2 g3，T_D = i g^1 g^3。
% ============================================================================
\section{过程依赖：Sivers 函数的变号}
\label{sec:sivers}

按 Trento 约定\cite{BacchettaEtAl2004}，横向极化核子（$\hat{\boldsymbol P}=\hat z$）中
\begin{equation}
  f_{q/p^\uparrow}(x,\boldsymbol k_{\mathrm T})=f_1(x,\kT^2)-f_{1T}^\perp(x,\kT^2)\,\frac{(\hat{\boldsymbol P}\times\boldsymbol k_{\mathrm T})\cdot\boldsymbol S}{M}
  =f_1-f_{1T}^\perp\frac{\epsilon^{ij}k_{\mathrm T}^iS_{\mathrm T}^j}{M},\qquad\epsilon^{12}=+1,
  \label{eq:sivers-def}
\end{equation}
它由 \eqref{eq:tmd-unsub} 中依赖自旋的部分给出；$[+]$（SIDIS）与 $[-]$（DY）分别对应未来与过去指向的订书钉。

\begin{theorem}[Sivers 函数变号]\tagR
\label{thm:sivers-sign}
对算符定义的 TMD（\cref{def:tmd-unsub,def:tmd-collins}），$\tilde f^{[+]}(x,\boldsymbol b_{\mathrm T};S)=\tilde f^{[-]}(x,\boldsymbol b_{\mathrm T};-S)$，因而
\begin{equation}
  f_1^{\mathrm{SIDIS}}(x,\kT^2)=f_1^{\mathrm{DY}}(x,\kT^2),\qquad f_{1T}^{\perp\,\mathrm{SIDIS}}(x,\kT^2)=-f_{1T}^{\perp\,\mathrm{DY}}(x,\kT^2).
  \label{eq:sivers-sign}
\end{equation}
\end{theorem}
\begin{proof}
取反幺正算符 $\mathcal A=\mathcal{PT}$。
\begin{enumerate}[label=(\roman*),leftmargin=2.2em]
\item 场：$\mathcal T\psi(x)\mathcal T^{-1}=T_D\psi(\tilde x)$，$\tilde x=(-x^0,\boldsymbol x)$，$\mathcal P\psi(x)\mathcal P^{-1}=\gamma^0\psi(\bar x)$；在 Dirac 表象中
$T_D=\ii\gamma^1\gamma^3$ 满足 $T_D\gamma^{\mu*}T_D^{-1}=\gamma_\mu$、$\gamma^0T_D^\dagger\gamma^0=T_D$、$T_D^2=1$。反幺正性使 c 数矩阵取复共轭，故
$\mathcal T\,\bar\psi(x)\Gamma\psi(y)\,\mathcal T^{-1}=\bar\psi(\tilde x)\,\gamma^0T_D^\dagger\gamma^0\,\Gamma^*\,T_D\,\psi(\tilde y)$；于是
$\gamma^+\xrightarrow{\ \mathcal T\ }\gamma^-\xrightarrow{\ \mathcal P\ }\gamma^+$。（单独的 $\mathcal T$ 把 $+z$ 方向强子的 TMD 变成 $-z$ 方向、含 $\gamma^-$ 的 TMD，
$\mathcal P$ 或绕横轴转 $\pi$ 再把它变回。）
\item 规范场：$\mathcal{PT}A^\mu(x)(\mathcal{PT})^{-1}=A^\mu(-x)$，色矩阵的复共轭由 $\eta_at^{a*}=t^a$ 抵消。对 Wilson 线，$\ii\to-\ii$ 与
$\dd z\to-\dd z$ 相互抵消，且反幺正变换保持算符乘积次序，故 $\mathcal{PT}\,W[b,a]\,(\mathcal{PT})^{-1}=W[-b,-a]$（沿反射路径、相同的路径编序）；
特别地 $\mathcal W^{[+]}(w,0;n)\to\mathcal W^{[-]}(-w,0;n)$，横向连接一并变换。
\item 态：$\mathcal{PT}|P,S\rangle=\eta|P,-S\rangle$，即动量不变、自旋矢量反号。用旋量可直接验证：
$\gamma^0T_D\,u^*(P,S)$ 所对应的投影算符正比于 $(\sla P+M)(1-\gamma_5\sla S)/2$。
\item 由 $\langle\alpha|O|\alpha\rangle=\langle\mathcal A\alpha|\mathcal AO\mathcal A^{-1}|\mathcal A\alpha\rangle^*$，记 $O^{[\pm]}(w)=\bar\psi(w)\gamma^+\mathcal W^{[\pm]}(w,0)\psi(0)$，
\[
  \begin{aligned}
  \langle P,S|O^{[+]}(w)|P,S\rangle&=\langle P,-S|O^{[-]}(-w)|P,-S\rangle^*=\langle P,-S|\big[O^{[-]}(-w)\big]^\dagger|P,-S\rangle\\
  &=\langle P,-S|O^{[-]}(w)|P,-S\rangle ,
  \end{aligned}
\]
最后一步用 $[O^{[-]}(-w)]^\dagger=\bar\psi(0)\gamma^+\mathcal W^{[-]}(0,-w)\psi(-w)$ 与平移 $+w$。对 $w^-$ 作 Fourier 变换即得
$\tilde f^{[+]}(S)=\tilde f^{[-]}(-S)$；Sivers 结构对 $S$ 为奇、$f_1$ 为偶。
\end{enumerate}
Collins\cite{Collins2002} 最先给出此论证。
\end{proof}

\begin{remark}[地位与适用条件]
\label{rem:sivers-status}
定理只用到 QCD 的 P、T 不变性、厄米性与平移不变性，对算符定义的函数是严格的（$R_{\mathrm{def}}$）；把它翻译成 SIDIS 与 DY 中
单自旋不对称性的反号则需要 TMD 因子化 \tagA。证明要求订书钉的路径编序与产生它的因子化推导中的编序一致：对类光或类空
（Collins）Wilson 线这自动成立；对\emph{类时}过去指向线，路径编序是反时间编序，而 Feynman 规则给出的是时间编序，二者在有限快度下
相差一个 Coulomb 相位（\cref{rem:timelike-lines}），此时 PT 像并不等于 DY 的 TMD。
\end{remark}

\paragraph{$b$ 空间形式。}由 $\int\dd^2k\,\ee^{-\ii\boldsymbol k\cdot\boldsymbol b}k^ig(k^2)=2\ii b^i\,\partial\tilde g/\partial b^2$（CONTEXT 的 Fourier 约定），
\begin{equation}
  \begin{gathered}
  \tilde f(x,\boldsymbol b_{\mathrm T};S)=\tilde f_1(x,b_{\mathrm T})+\ii M\,\epsilon^{ij}b_{\mathrm T}^iS_{\mathrm T}^j\,\tilde f_{1T}^{\perp(1)}(x,b_{\mathrm T}),\\
  \tilde f_{1T}^{\perp(1)}\equiv-\frac{2}{M^2}\frac{\partial\tilde f_{1T}^\perp}{\partial b_{\mathrm T}^2},\qquad
  \tilde f_{1T}^{\perp(1)}(x,0)=\int\dd^2k_{\mathrm T}\frac{\kT^2}{2M^2}f_{1T}^\perp ,
  \end{gathered}
  \label{eq:sivers-bspace}
\end{equation}
Sivers 项为纯虚，符合 $\tilde f(x,\boldsymbol b)^*=\tilde f(x,-\boldsymbol b)$；定理即 $\tilde f_{1T}^{\perp(1)[+]}=-\tilde f_{1T}^{\perp(1)[-]}$。

\begin{example}[单胶子末态相互作用的相位：标量双夸克模型]
\label{ex:sivers-model}
核子（质量 $M$）以点耦合 $g$ 分裂为夸克（$m$）与标量双夸克（$\lambda$）；阿贝尔胶子与夸克（及其 Wilson 线）耦合 $e_q$、与双夸克耦合 $e_s$。
树图 $\mathcal M_0=-g\frac{\sla k+m}{k^2-m^2}U(P,S)$，$k^2-m^2=-\frac{\kT^2+L^2}{1-x}$，$L^2=x\lambda^2+(1-x)m^2-x(1-x)M^2$。
Wilson 线与末态双夸克之间交换一个胶子：
\begin{equation}
  \begin{gathered}
  \mathcal M_1=-\ii e_qe_sg\!\int\!\frac{\dd^4l}{(2\pi)^4}\frac{[2(1-x)P^++l^+]\,(\sla k-\sla l+m)\,U}{(l^+\mp\ii0)\,\mathcal D(l)}\qquad(\text{SIDIS 取}-,\ \text{DY 取}+),\\
  \mathcal D(l)=(l^2+\ii0)\big((P-k+l)^2-\lambda^2+\ii0\big)\big((k-l)^2-m^2+\ii0\big).
  \end{gathered}
  \label{eq:sivers-M1}
\end{equation}
T-奇部分来自 $\frac{1}{l^+\mp\ii0}=\mathrm{PV}\frac1{l^+}\pm\ii\pi\delta(l^+)$。在 $l^+=0$ 处用留数做 $l^-$ 积分
（$=\frac{\pi\ii}{P^+[(\boldsymbol k-\boldsymbol l)_{\mathrm T}^2+L^2]}$），$l^-\gamma^+$ 项因 $\gamma^+\gamma^+=0$ 消失：
\begin{equation}
  \begin{gathered}
  \mathcal M_1^{\mathrm{odd}}=\mp2\pi^2\ii(1-x)\frac{e_qe_sg}{(2\pi)^4}\int\dd^2l_{\mathrm T}\,
  \frac{(\sla k+\boldsymbol\gamma_{\mathrm T}\cdot\boldsymbol l_{\mathrm T}+m)\,U}{\boldsymbol l_{\mathrm T}^2\,[(\boldsymbol k-\boldsymbol l)_{\mathrm T}^2+L^2]},\\
  \int\dd^2l_{\mathrm T}\frac{\boldsymbol l_{\mathrm T}}{\boldsymbol l_{\mathrm T}^2[(\boldsymbol k-\boldsymbol l)_{\mathrm T}^2+L^2]}=\frac{\pi\boldsymbol k_{\mathrm T}}{\kT^2}\ln\frac{\kT^2+L^2}{L^2}.
  \end{gathered}
  \label{eq:sivers-M1odd}
\end{equation}
以 $f=\frac{1}{4(2\pi)^3(1-x)P^+}\sum\overline{\mathcal M}\gamma^+\mathcal M$、$U\bar U=(\sla P+M)(1+\gamma_5\sla S)/2$ 得
\begin{equation}
  \begin{gathered}
  f_1=\frac{g^2}{16\pi^3}\frac{(1-x)[(m+xM)^2+\kT^2]}{(\kT^2+L^2)^2},\\
  f_{1T}^{\perp\,\mathrm{SIDIS}}=\frac{g^2e_qe_s}{4(2\pi)^4}\frac{M(m+xM)(1-x)}{\kT^2(\kT^2+L^2)}\ln\frac{\kT^2+L^2}{L^2}=-f_{1T}^{\perp\,\mathrm{DY}}.
  \end{gathered}
  \label{eq:sivers-model}
\end{equation}
吸引的末态相互作用（$e_qe_s<0$；QCD 色单态 $e_qe_s\to-4\pi C_F\as$）给出 $f_{1T}^{\perp\,\mathrm{SIDIS}}<0$（$m+xM>0$）。这与
Bacchetta--Conti--Radici\cite{BacchettaContiRadici2008} 点耦合标量双夸克的结果完全一致（其程函顶角 $-\ii e_cn^\rho$、传播子 $\frac{\ii}{-l^++\ii\epsilon}$、
双夸克顶角 $+\ii e_c(2P-2p+l)^\rho$ 对应本书的 $e_qe_s=-e_c^2$，且 $e_c^2\to4\pi C_F\as$）。物理图像（Brodsky--Hwang--Schmidt\cite{BrodskyHwangSchmidt2002}）：
被击中夸克与旁观者之间的在壳（类 Coulomb）再散射相位，与树图振幅在相差一个单位轨道角动量的核子螺旋度态之间干涉（因子 $m+xM$）；
SIDIS 的末态相互作用与 DY 的初态相互作用的程函传播子 $\ii0$ 反号，因而不对称性反号。
\end{example}
